\section{Background} \label{sec:background}

\paragraph{Boolean Circuits and Boolean Networks} Boolean circuits are 
directed acyclic graphs where each node computes a bivariate 
Boolean function $B:\{0,1\}^2\rightarrow \{0,1\}$ \cite{arora2009computational}.
We denote the $16$ possible bivariate Boolean functions by 
$\sB:= \{B_i: \{0,1\}^2 \rightarrow \{0,1\}, i=1,\ldots, 16\}$ (c.f. \cref{app:bool_fun}).
A Boolean network is a Boolean circuit 
where nodes and edges are arranged in layers. \cref{fig:boolean_circuit_and_network} illustrates a Boolean circuit and a Boolean network.

\begin{figure}
\centering
\resizebox{.7\linewidth}{!}{\begin{tikzpicture}[
  gate/.style={circle, draw, minimum size=4mm, inner sep=0pt, font=\small},
  io/.style={font=\small},
  wire/.style={-}
]

\node[io] (a) at (1,0) {$x_1$};
\node[io] (b) at (1,1) {$x_2$};
\node[io] (c) at (1,2) {$x_3$};
\node[io] (d) at (6,0.5) {};
\node[io] (e) at (6,1.5) {};
\node[gate] (g1) at (2,2) {$B_1$};
\node[gate] (g2) at (3,1) {$B_4$};
\node[gate] (g3) at (5,0.5) {$B_7$};
\node[gate] (g4) at (5,1.5) {$B_3$};



\draw[-Stealth] (c.east) -- (g1.west);
\draw[-Stealth] (b.east) -- (g1.west);
\draw[-Stealth] (g1.east) -- (g2.west);
\draw[-Stealth] (a.east) -- (g2.west);
\draw[-Stealth] (a.east) -- (g3.west);
\draw[-Stealth] (g1.east) -- (g3.west);
\draw[-Stealth] (g3.east) -- (d.west);

\draw[-Stealth] (g2.east) -- (g4.west);
\draw[-Stealth] (g1.east) -- (g4.west);

\draw[-Stealth] (g4.east) -- (e.west);


\begin{scope}[xshift=7cm]

\node[io] (a) at (1,0) {$x_1$};
\node[io] (b) at (1,1) {$x_2$};
\node[io] (c) at (1,2) {$x_3$};
\node[io] (o1) at (4,1) {};
\node[io] (o2) at (4,2) {};
\node[io] (o3) at (4,0) {};

\node[gate] (g11) at (2,1) {$B_1$};
\node[gate] (g12) at (2,2) {$B_2$};
\node[gate] (g13) at (2,0) {$B_5$};
\node[gate] (g21) at (3,1) {$B_7$};
\node[gate] (g22) at (3,2) {$B_2$};
\node[gate] (g23) at (3,0) {$B_4$};

\draw[-Stealth] (a.east) -- (g11.west);
\draw[-Stealth] (b.east) -- (g11.west);
\draw[-Stealth] (b.east) -- (g12.west);
\draw[-Stealth] (c.east) -- (g12.west);
\draw[-Stealth] (c.east) -- (g13.west);
\draw[-Stealth] (a.east) -- (g13.west);


\draw[-Stealth] (g12.east) -- (g21.west);
\draw[-Stealth] (g13.east) -- (g21.west);
\draw[-Stealth] (g12.east) -- (g22.west);
\draw[-Stealth] (g13.east) -- (g22.west);
\draw[-Stealth] (g11.east) -- (g23.west);
\draw[-Stealth] (g13.east) -- (g23.west);
\draw[-Stealth] (g23.east) -- (o3.west);
\draw[-Stealth] (g22.east) -- (o2.west);
\draw[-Stealth] (g21.east) -- (o1.west);


\end{scope}
\end{tikzpicture}
}
\caption{Left: a Boolean circuit. Right: a Boolean network.}\label{fig:boolean_circuit_and_network}
\vspace{-1em}
\end{figure}

Boolean circuits realize only Boolean functions.
To apply them to implement a non-Boolean function $g:A\rightarrow B$, 
we need to decompose $g$ into a Boolean function $f:\{0,1\}^{m} \rightarrow \{0,1\}^n, m,n\in \sN$, together with 
an encoder $\gE:A \rightarrow \{0,1\}^m$ and a decoder $\gD: \{0,1\}^n\rightarrow B$, such that $g=\gD  \circ f\circ \gE$.
In particular, for an image classification function 
  $g:[0,1]^d \rightarrow \gC$  that takes a real-valued image and maps it to a class label, we %
  use the thermometer encoder, defined as $\gE(\vx) := [\mathbbm{1}\{\vx \ge \frac{1}{N}\}, \dots, \mathbbm{1}\{\vx \ge \frac{N-1}{N}\}]$, where is $\mathbbm{1}\{\cdot\}$ the indicator function. We use the population count decoder (also named GroupSum) selecting the class with
 the highest number of 1s in the output Boolean vector $\mY \in \{0,1\}^{|\gC| \times n_c}$, defined as $\gD(\mY):= \argmax_{c \in \gC} \sum_{i=1}^{n_c} \mY_{c,i}$.


\paragraph{Learning Boolean Networks via Differentiable Relaxations} To learn a Boolean network from data, \citet{petersen2022deep} propose to 
relax  discrete Boolean networks into differentiable functions, enabling training via gradient descent.
Specifically, for $i=1,\ldots, 16$, let $B_i^c:[0,1]^2\rightarrow [0,1]$ be a differentiable relaxation of $B_i$. 
For example, $B_2(x_b,y_b) = x_b\wedge y_b, x_b,y_b\in \{0,1\}$, is relaxed into $B_2^c(x,y) = xy, x,y\in [0,1]$, satisfying $B_2^c(x_b, y_b) = B_2(x_b, y_b)$.
Each neuron is then parametrized as a weighted aggregation of all $B_i^c$ as 
\begin{equation} \label{eq:relaxation}
    f^c(x,y) = \sum_{i=1}^{16} \frac{\exp(\vw_i)}{\sum_{j=1}^{16} \exp(\vw_j)} B_i^c(x,y),
\end{equation}
where $\vw_i$ are the parameters.
Afterwards, the relaxed network is trained on data by updating the weight vector $\vw = (\vw_1,\ldots,\vw_{16})$ of each neuron. 
Finally, each neuron is discretized back to a bivariate Boolean function by 
\begin{equation} \label{eq:discretization}
    f := B_{i^*},  \quad \text{ with } i^* = \argmax_{i} \vw_i.
\end{equation}
Note that the relaxation and parametrization above only allows 
to learn the Boolean operation at each neuron, but does not allow the learning of network connections. 
We shall address the efficient connection learning  later in \cref{sec:method}.

\paragraph{Convolutional Boolean Networks} 
Convolutional layers are key components in floating-point neural networks 
capturing local spatial structures. 
\citet{petersen2024convolutional} design a convolutional architecture for Boolean networks,  where each kernel is a binary tree consisting of Boolean operations. 
A tree of depth $d$ requires $2^d-1$ Boolean operations to 
implement a Boolean function with $2^{d-1}$ inputs. As a result, their convolutional layer 
costs significantly more Boolean operations than a non-convolutional layer. 


